\documentclass[a4paper,11pt]{article}
\usepackage[T1]{fontenc}
\usepackage[utf8]{inputenc}
\usepackage[italian]{babel}
\usepackage{amsmath,amssymb}
\usepackage{graphicx}
\usepackage{geometry}
\usepackage{xcolor}
\usepackage{listings}
\geometry{margin=2.5cm}

\lstdefinelanguage{MiniZinc}{
  morekeywords={int,set,of,array,var,constraint,include,solve,satisfy,output,show},
  sensitive=true,
  morecomment=[l]{\%},
  morestring=[b]"
}

\lstset{
  language=MiniZinc,
  basicstyle=\ttfamily\small,
  keywordstyle=\color{blue}\bfseries,
  commentstyle=\color{teal!70!black}\itshape,
  stringstyle=\color{orange!70!black},
  numbers=left,
  numberstyle=\tiny\color{gray},
  stepnumber=1,
  numbersep=8pt,
  showstringspaces=false,
  frame=single,
  breaklines=true
}

\title{CSP: Cards}
\date{}

\begin{document}
\maketitle

È data una collezione di carte, in cui ad ogni carta è associato un intero.
La collezione è costituita da $k > 2$ carte per ogni intero da $1$ a $n$, con $n > 0$.
Vogliamo disporre in fila le $k \cdot n$ carte in modo che ogni carta contrassegnata dal numero $i$ sia in posizione tale da avere esattamente $i$ carte che la dividono dalla precedente carta contrassegnata da $i$ (se esiste).


\subsection*{1.1 Modellazione} 
Una formulazione CSP ad alto livello è la seguente.

\paragraph{Variabili}
Definiamo una variabile per ogni carta:
\[
X = \{x_{ij} \mid i \in \{1, \ldots, n\}, j \in \{1, \ldots, k\}\},
\]
Dove $i$ rappresenta l'intero associato alla carta e $j$ rappresenta l'occorrenza della carta con quel numero.
Il valore $x_{ij}$ rappresenta la posizione della carta.
\paragraph{Domini}
\[
\forall x \in X:\quad D(x)=\{1,\ldots,n*k\}.
\]

\paragraph{Vincoli}

Ogni carta occupa una posizione unica:
\[
C_1 = \text{alldifferent}(\{x_{ij} \mid i \in \{1, \ldots, n\}, j \in \{1, \ldots, k\}\}).
\]

Per ogni carta che non sia la prima occorrenza di un numero, la carta deve essere posizionata a una posizione pari alla precedente occorrenza di quel numero più $i+1$ (dove $i$ è il numero associato alla carta), in modo da garantire una distanza di $i$ carte tra esse:
\[
C_2 = \{\forall i \in \{1, \ldots, n\}, j \in \{2, \ldots, k\}:\quad x_{ij} = x_{i(j-1)} + (i+1)\}.
\]


\[
C = C1 \land C2
\]

\newpage

\subsection*{1.2 Istanziazione per (k, n) = (2, 4)}


Si ottiene il CSP $(X,D,C)$:

\[
X = \{
x_{1,1}, x_{1,2}, x_{2,1}, x_{2,2}, x_{3,1}, x_{3,2}, x_{4,1}, x_{4,2}
\}
\]

\[
\forall x \in X:\quad D(x)=\{1, \ldots, 8\}.
\]

\[
C_1 = \text{alldifferent}(\{x_{1,1}, \ldots, x_{4,2}\})
\]

\[
C_2 = \left\{
\begin{aligned}
&x_{1,2} = x_{1,1} + (1+1),\\
&x_{2,2} = x_{2,1} + (2+1),\\
&x_{3,2} = x_{3,1} + (3+1),\\
&x_{4,2} = x_{4,1} + (4+1)
\end{aligned}
\right\}
\]

\[
C = C_1 \land C_2
\]

\newpage

\subsection*{1.3 Codifica in MiniZinc}
\begin{lstlisting}
include "alldifferent.mzn";
int: n;
int: k;
array[1..n, 1..k] of var 1..n*k: x;

constraint alldifferent(x);
constraint forall(i in 1..n, j in 2..k) (
    x[i,j] = x[i,j-1] + (i+1)
);

solve satisfy;

output [show(x)];
\end{lstlisting}

Si può visualizzare la soluzione con la sequenza di carte con questo codice:

\begin{lstlisting}
include "alldifferent.mzn";
int: n;
int: k;
array[1..n, 1..k] of var 1..n*k: x;

constraint alldifferent(x);
constraint forall(i in 1..n, j in 2..k) (
    x[i,j] = x[i,j-1] + (i+1)
);
 

% trasformazione in sequenza
array[1..k*n] of var 1..n: sequence;
constraint forall(i in 1..n, j in 1..k)(
    sequence[x[i, j]] = i
);

solve satisfy;

output [show(sequence)];
\end{lstlisting}



\end{document}